% answers_short.tex —— 2014 年·辽宁卷·理科 · 简短答案速查
% 由 tools/make_short_answers.py 自 answers.tex 的 \gksolhead/\gkanswer 提取，
% 并与 2014/2014_数学真题_辽宁理_答案要点.md 逐题比对后生成，勿手改（改源 MD 后重跑）。

\begin{gknotice}
\textbf{使用说明}：本节只给最终答案便于快速核对，完整推导见解析版。答案为编写者独立推导，\textbf{非官方标准答案}。
\end{gknotice}

\gkpart{选择题}
\begin{center}\begin{tabular}{ccccccccccccc}
\toprule
题号 & 1 & 2 & 3 & 4 & 5 & 6 & 7 & 8 & 9 & 10 & 11 & 12 \\
\midrule
答案 & D & A & C & B & A & D & B & D & B & D & C & B \\
\bottomrule
\end{tabular}\end{center}

\gkpart{填空题}
\begin{itemize}[leftmargin=2.6em]
  \item[\textbf{13.}] $\dfrac{29}9$。依次取 $x=9,5,\dfrac{11}3$，第三次得 $y=\dfrac{29}9$，此时 $|y-x|=\dfrac49<1$，输出 $\dfrac{29}9$。
  \item[\textbf{14.}] $\dfrac13$。阴影面积 $\displaystyle\int_{-1}^{1}2x^2\,\mathrm dx=\dfrac43$，正方形面积 $4$，概率为 $\dfrac13$。
  \item[\textbf{15.}] $12$。设 $Q$ 为 $MN$ 中点，则 $|AN|=2|F_1Q|$，$|BN|=2|F_2Q|$，故和为 $2\times2a=12$。
  \item[\textbf{16.}] $-2$。由 $4a^2-2ab+4b^2\geqslant\dfrac58(2a+b)^2$ 知取最大值时 $b=\dfrac23a$、$c=\dfrac{40}9a^2$，原式化为 $\dfrac5c-\dfrac{2\sqrt{10}}{\sqrt c}$，当 $c=\dfrac52$ 时取最小值 $-2$。
\end{itemize}

\gkpart{解答题}
\begin{itemize}[leftmargin=2.6em]
  \item[\textbf{17.}] (1) $a=3$，$c=2$。由 $ac\cos B=2$ 得 $ac=6$，再由余弦定理得 $a^2+c^2=13$，结合 $a>c$ 解出；(2) $\dfrac{23}{27}$。$\cos C=\dfrac79$，$\sin C=\dfrac{4\sqrt2}9$，$\sin B=\dfrac{2\sqrt2}3$，代入两角差余弦公式
  \item[\textbf{18.}] (1) $0.108$。日销售量不低于 $100$ 的概率为 $0.6$，低于 $50$ 的概率为 $0.15$，所求为 $2\times0.6^2\times0.15=0.108$；(2) $P(X=0)=0.064$，$P(X=1)=0.288$，$P(X=2)=0.432$，$P(X=3)=0.216$；$E(X)=1.8$，$D(X)=0.72$。$X\sim B(3,0.6)$，直接代入二项分布公式
  \item[\textbf{19.}] (1) 见解析版；(2) $\dfrac{2\sqrt5}5$。平面 $BEF$ 与平面 $BCF$ 的法向量分别为 $(3,-\sqrt3,-\sqrt3)$ 与 $(0,1,0)$，二面角余弦的绝对值为 $\dfrac1{\sqrt5}$，故正弦为 $\dfrac{2\sqrt5}5$
  \item[\textbf{20.}] (1) $x^2-\dfrac{y^2}2=1$。由 $e=\sqrt3$ 得 $b^2=2a^2$，代入 $P(\sqrt2,\sqrt2)$ 得 $a^2=1$；(2) $y=\dfrac{3\sqrt6+2}{25}(x-\sqrt3)$ 或 $y=-(2+\sqrt6)(x-\sqrt3)$。设 $l:y=k(x-\sqrt3)$，由 $\overrightarrow{PA}\cdot\overrightarrow{PB}=0$ 得 $(11-4\sqrt6)k^2+2\sqrt6k-2=0$，解出 $k$
  \item[\textbf{21.}] (1) 见解析版；(2) 见解析版
\end{itemize}

\gkpart{选做题}
\begin{itemize}[leftmargin=2.6em]
  \item[\textbf{22.}] (1) 见解析版；(2) 见解析版
  \item[\textbf{23.}] (1) $\begin{cases}x=\cos\theta,\\y=2\sin\theta\end{cases}$（$\theta$ 为参数）。由 $x^2+\dfrac{y^2}4=1$ 直接写出；(2) $\rho=\dfrac3{4\sin\theta-2\cos\theta}$。中点 $M\left(\dfrac12,1\right)$，所求直线 $2x-4y+3=0$，代入 $x=\rho\cos\theta$，$y=\rho\sin\theta$
  \item[\textbf{24.}] (1) $M=\left[0,\dfrac43\right]$。分 $x\geqslant1$ 与 $x<1$ 去绝对值求解；(2) 见解析版
\end{itemize}
